\section*{Sequential 6-Bit Palindrome Detector Instructions}

\subsection*{Objective}
Design a sequential digital circuit that reads one input bit per clock cycle and determines whether each 6-bit word is a palindrome.

\subsection*{Allowed Components}
Use all basic logic gates, plus DFF for state:
\begin{itemize}
  \item AND
  \item OR
  \item NOT
  \item XOR
  \item NAND
  \item NOR
  \item XNOR
  \item DFF
\end{itemize}

\subsection*{Inputs and Outputs}
\begin{itemize}
  \item Input: $X$ (serial bit stream, one bit each cycle)
  \item Output: is_palindrome
\end{itemize}

\subsection*{Timing Model}
Each test stream contains exactly 6 cycles and represents one word:
\[
b1,b2,b3,b4,b5,b6
\]
The circuit is reset before each stream (all DFF states start at 0).

\subsection*{What Is a 6-Bit Palindrome?}
A 6-bit word is a palindrome if:
\[
b1=b6, & b2=b5, & b3=b4
\]


\subsection*{Output Convention}
\begin{itemize}
  \item For cycles 1--5, output is_palindrome=0
  \item On cycle 6, output 1 if all mirrored pairs matched, otherwise 0
\end{itemize}

\subsection*{Examples}
\begin{itemize}
  \item Input stream: \texttt{1,0,1,1,0,1} (word \texttt{101101}) \\
  Expected output stream: \texttt{0,0,0,0,0,1}
  \item Input stream: \texttt{1,1,0,1,0,0} (word \texttt{110100}) \\
  Expected output stream: \texttt{0,0,0,0,0,0}
\end{itemize}

\subsection*{Hint}
A common design is to store the first three bits using DFFs, accumulate any mismatch in a state bit, and output the negation of that mismatch bit on cycle 6.